\chapter{File Handling \& Exception Management}

Data science datasets rarely start in memory—they live on disk as CSV, JSON, text logs, or database dumps. In this chapter, we cover \textbf{File I/O} (Input/Output) operations and \textbf{Exception Handling} to ensure robust data ingestion pipelines that don't crash when encountering missing or corrupted data.

\section{Reading and Writing Files}

Python uses the built-in `open(filename, mode)` function to interact with files on disk. Using the Context Manager `with` guarantees that file handles are automatically closed, even if errors occur.

\begin{definitionbox}{File Modes and Context Managers}
Common modes: `'r'` (read), `'w'` (write / overwrite), `'a'` (append), `'b'` (binary mode).
\begin{lstlisting}
# Writing to a text file
with open("data.txt", "w") as f:
    f.write("Line 1: Data Science\n")
    f.write("Line 2: Machine Learning\n")

# Reading line by line (memory efficient)
with open("data.txt", "r") as f:
    for line in f:
        print(line.strip())
\end{lstlisting}
\end{definitionbox}

\begin{videocard}{IITM BS Lecture: Reading and Writing to a File}{https://www.youtube.com/watch?v=rYLJaAdgLhI}
Official lecture demonstrating file modes, `open()`, `.read()`, `.readline()`, and `with` context blocks.
\end{videocard}

\section{Processing Large Datasets}

When log files or text datasets exceed available RAM, loading the entire file with `.read()` causes an Out-Of-Memory (OOM) crash. Instead, stream the file line by line using an iterator or chunked reads.

\begin{theorembox}{Line-by-Line Streaming}
Iterating over a file object `for line in f:` streams one line at a time into memory ($\mathcal{O}(1)$ memory complexity).
\end{theorembox}

\begin{videocard}{IITM BS Lecture: Handling Very Big Files}{https://www.youtube.com/watch?v=NwFpo_KGvOA}
Official lecture discussing memory considerations when searching and processing massive files.
\end{videocard}

\section{Exception Handling (`try`, `except`, `else`, `finally`)}

An \textbf{Exception} is a runtime error that halts program execution. Exception handling allows Python to catch runtime errors gracefully.

\begin{formulabox}{The Exception Handling Block}
\begin{lstlisting}
try:
    # Code that might raise an exception
    num = int(input("Enter a number: "))
    result = 100 / num
except ValueError:
    print("Invalid input! Please enter a number.")
except ZeroDivisionError:
    print("Division by zero is undefined.")
else:
    # Executed ONLY if NO exception occurred
    print(f"Result: {result}")
finally:
    # Executed ALWAYS (used for cleanup)
    print("Execution complete.")
\end{lstlisting}
\end{formulabox}

\textbf{Data Science Relevance:} Web scrapers, data loaders, and model training loops use `try-except` blocks to skip corrupted rows or bad images without crashing overnight training runs.

\begin{videocard}{IITM BS Lecture: Exception Handling Explained}{https://www.youtube.com/watch?v=me3_XR7Iz60}
Official lecture explaining `try`, `except`, `finally`, built-in exception types, and raising custom errors.
\end{videocard}

\newpage
\section{Practice Exercises}
\textit{Detailed step-by-step solutions are available in the Appendix.}

\begin{enumerate}
    \item \textbf{File Writing \& Reading:} Write a Python script that:
    \begin{enumerate}
        \item Creates a file named `log.txt` and writes $3$ lines of sample log entries.
        \item Reopens `log.txt` in read mode and counts the total number of lines in the file.
    \end{enumerate}

    \item \textbf{Handling Missing Files:} Write a script using `try-except` that attempts to open `"missing_dataset.csv"`. If the file does not exist, catch `FileNotFoundError` and print `"Dataset file not found. Please check filepath."`.

    \item \textbf{Safe Type Conversion Function:} Write a function `safe_float(value)` that attempts to convert `value` to a `float`. If a `ValueError` or `TypeError` occurs, return `None`. Test with `"3.14"`, `"abc"`, and `[1, 2]`.

    \item \textbf{Parsing Log File Data:} Suppose a text file `scores.txt` contains numeric test scores, one per line. Write a script to read the file, convert each line to float inside a `try` block (skipping invalid lines), and calculate the average score.

    \item \textbf{Data Science Application (CSV Line Counter & Cleaner):} Write a script `clean_csv(input_path, output_path)` that reads a CSV file line by line, removes empty lines or lines starting with `#` (comments), and writes clean lines to `output_path`.
\end{enumerate}